%! Author = Saeid Yazdani
%! Date = 12/3/2023


\section{The Mixed-Signal Era}
\label{sec:mixed-signal-era}

The invention of the transistor in 1947 and the Integrated Circuit in 1958
established a path for complex chip designs.
In the beginning, analog and digital circuits were developed separately, with
analog solutions being dominant in almost all electronic systems.

In the 1990s, the semiconductor industry made a tremendous change of
direction regarding chip design and test philosophy: Integrating analog and
digital function blocks on the same \ac{IC} die and package, referred to as
a \emph{mixed-signal} \ac{IC}.
Mixed-signal chips are involved in interfacing analog signals like a
sensor or an \ac{RF} antenna with digital processing power in one package.
A Mixed-signal \ac{IC} typically combines analog functions like
an \ac{ADC} or operational amplifiers with digital functions such as
\ac{DSP} cores.

In earlier days, there was a clear partitioning of \ac{IC}s into analog,
digital, memory and \ac{RF}.
For this very reason, the \acp{ATE} were strictly partitioned into separate
categories of \emph{Analog}, \emph{Logic (digital)}, \emph{Memory}
and \emph{\acs{RF}} testers.

In the 1990s, the interest in the development of mixed-signal \acp{IC}
began to rise, and these chips started to gain
significant market share.
Mixed-signal \acp{IC} found general use in diverse
applications from telecommunications to consumer electronics.

The development of \ac{CMOS} technology was particularly significant, as it allowed
the integration of high-density digital logic with analog circuits on
the same chip.
At this time:

\begin{itemize}[topsep=8pt,itemsep=4pt,partopsep=4pt, parsep=4pt]
    \item Digital logic was still using large transistors
    operating at \SI{5}{\volt}
    \item Design technology succeeded in the development of amplifiers with the
    same technology (analog building blocks)
\end{itemize}

In the 2000s, a trend started toward smaller and more power-efficient
mixed-signal \ac{IC}s, as a response to the increased demand for portable
electronic devices.


This major revolution has raised the need for new test methodologies to support
this new type of integrated circuits: Standard \ac{ATE} systems were no longer
sufficient, and mixed-signal \ac{ATE} for \ac{SoC} were developed that operate
in both analog and digital domains in one mainframe \autocite{Tsuboshita}.

As can be observed in \autoref{tab:analog-vs-digital-signals}, there are
fundamental differences between analog and digital.
The complexity in mixed-signal testing arises from the need to accurately
evaluate the performance of both types of signals.

\begin{table}[]
    \resizebox{\textwidth}{!}{% Limit table size to text with by reducing font
        \centering
        \begin{tabular}{@{}lll@{}}
            \toprule
            Property             & Analog Signals                  & Digital Signals                    \\
            \midrule
            Data Flow            & Continuous                      & Discrete                           \\
            Useful Data          & Amplitude, period and frequency & High or Low                        \\
            Noise Susceptibility & Prone to noise and distortion   & Less prone to noise and distortion \\
            \bottomrule
        \end{tabular}
    } % End resize-box
    \caption[Analog v.s. Digital signals]{A list of major differences between
    analog and digital signals.}
    \label{tab:analog-vs-digital-signals}
\end{table}



A mixed-signal chip will operate in both digital and analog domains at the same
time to deliver a desired output.
Because analog and digital functions of a mixed-signal \ac{DUT} are not
independent, a mixed-signal \ac{ATE} contains compartments for digital logic
test (i.e. scan test), memory test, and analog block test.

For digital \ac{IC}s, \ac{ATE} focuses on logic functionality, timing, and
data integrity and the test itself involves applying digital test
patterns and comparing the output against pre-defined lists of correct values
(expected responses).
Techniques like \ac{ATPG}, SCAN and \ac{BIST} are well established.
The fact that digital signals are discrete in time and level makes it easier to
feed a \ac{DUT} with a set of test patterns.

On the contrary, testing analog \acp{IC} requires more complex
signal generation and analysis as the testing revolves around parameters
like voltage levels, frequency response, noise, and distortion.

Testing a mixed-signal chip that combines digital
and analog signals can prove difficult for multiple reasons, of which:

\begin{itemize}[topsep=8pt,itemsep=4pt,partopsep=4pt, parsep=4pt]
    \item Analog test signals must be synchronized to the ATE master clock
    (synchronous signal generation and synchronous signal acquisition)
    \item Analog signals exhibit static characteristics such as offset voltage
    or bias currents, and dynamic characteristics like bandwidth, gain, and
    phase response.
\end{itemize}

Testing costs time and money, and we want a \emph{single insertion test}
approach, therefore, it is desirable to test mixed-signal chips on a single
\ac{ATE}.
A \emph{one-size-fit-all} \ac{ATE}~\autocite{Burns} is sought
in the semiconductor industry.
The \ac{ATE} must be capable of delivering
and measuring analog and digital signals with high precision and
accurate timing.
\autoref{fig:one-size-fit-all-ate} Shows an abstract view of a mixed-signal
\ac{ATE}.

\begin{figure}[!h]
    \centering
    \printfitgraphics[0.36\textheight]{figures/analog-vs-digital-ate}
    \caption[Analog v.s. Digital ATE]{A Comparison between
    expected functionality and differences between analog and digital
    \ac{ATE}s. A Mixed-signal or \ac{SoC} enabled \ac{ATE} will have a
    combination of both analog and digital capabilities.}
    \label{fig:one-size-fit-all-ate}
\end{figure}

